% Part: incompleteness
% Chapter: theories-computability
% Section: oconsis-ext-of-q-undec

\documentclass[../../../include/open-logic-section]{subfiles}

\begin{document}

\olfileid{inc}{tcp}{oqn}

\olsection{$\Th{Q}$ యొక్క $\omega$-అవైరుధ్య విస్తరణలు అనిర్ణయనీయాలు}

\begin{explain}
$\Th{Q}$ గ.లె.-సంపూర్ణమనే నిరూపణ, $\Th{Q}$లో నిరూపించదగిన ప్రతి వాక్యమూ
సహజ సంఖ్యల విషయంలో ``సత్యం'' అనే వాస్తవంపై ఆధారపడింది. పై నిరూపణకు
అవసరమైన ``సత్యం'' యొక్క అంశాలను ఖచ్చితంగా గుర్తించడం ద్వారా, తదుపరి నిర్వచనం,
సిద్ధాంతం ఆ ఫలితాన్ని బలపరుస్తాయి. అయితే ఈ సిద్ధాంతంపై ఎక్కువసేపు
ఆగవద్దు; త్వరలోనే దీన్ని ఇంకా బలపరుస్తాం. దీన్ని ప్రధానంగా చారిత్రక కారణాల
కోసం చేర్చాం: గ్యోడెల్ తన అసలు వ్యాసంలో $\omega$-అవైరుధ్య భావనను వాడాడు;
కాని కొద్దికాలానికే $\omega$-అవైరుధ్య స్థానంలో సాధారణ అవైరుధ్యాన్ని పెట్టి
అతని ఫలితాన్ని బలపరిచారు.
\end{explain}

\begin{defn}
\ollabel{thm:oconsis-q}
కిందిది నెరవేరితే సిద్ధాంతం~$\Th{T}$ను $\omega$-అవైరుధ్యమైనది అంటాం:
$\lexists[x][!A(x)]$ ఏ వాక్యమైనా, $\Th{T}$ అనేది $\lnot
!A(\num 0)$, $\lnot !A(\num 1)$, $\lnot !A(\num 2)$, \dots లను
నిరూపిస్తే, $\Th{T}$ అనేది $\lexists[x][!A(x)]$ను నిరూపించదు.
\end{defn}

\begin{thm}
  $\Th{Q}$ను కలిగి ఉన్న ఏ $\omega$-అవైరుధ్య సిద్ధాంతమైనా $\Th{T}$గా
  తీసుకోండి. అప్పుడు $\Th{T}$ నిర్ణయించదగినది కాదు.
\end{thm}

\begin{proof}
$\Th{T}$, $\Th{Q}$ను కలిగి ఉంటే, గణనీయ ప్రమేయాలు, సంబంధాలకు $\Th{T}$
ప్రాతినిధ్యం వహిస్తుంది. ముందటి నిరూపణను కొద్దిగా మార్చితే చాలు. పైన చెప్పినట్లే,
$x \in K$ అయితే $\Th{T}$ అనేది $\lexists[s][!A_T(\num
  x,\num x, s)]$ను నిరూపిస్తుంది. తిరుగుగా, $\Th{T}$ అనేది
$\lexists[s][!A_T(\num x, \num x, s)]$ను నిరూపిస్తుందని అనుకుందాం. అప్పుడు
$x$, $K$లో ఉండాలి: లేకపోతే, ఇన్‌పుట్~$x$పై యంత్రం~$x$కు ఆగే గణన ఉండదు;
$!A_T$ క్లీనీ $T$ సంబంధానికి ప్రాతినిధ్యం వహిస్తుంది కనుక, $\Th{T}$ అనేది
$\lnot !A_T(\num x, \num x, \num 0)$, $\lnot !A_T(\num x, \num
x, \num 1)$, \dots లను నిరూపిస్తుంది. దీనివల్ల $\Th{T}$
$\omega$-వైరుధ్యమైనదవుతుంది.
\end{proof}

\end{document}
